% !TEX program = lualatex
% Transparencias de Fundamentos de las Matemáticas I
% Clase 13: Semigrupos, monoides y grupos

\documentclass[aspectratio=169,11pt]{beamer}

\usepackage{fontspec}
\usepackage[spanish,es-nodecimaldot]{babel}
\usepackage{amsmath,amssymb,mathtools}
\usepackage{booktabs}
\usepackage{csquotes}
\usepackage{xcolor}

\usetheme{Madrid}
\usecolortheme{default}
\setbeamertemplate{navigation symbols}{}
\setbeamertemplate{footline}[frame number]

\definecolor{FdMBlue}{RGB}{20,55,100}
\definecolor{FdMRed}{RGB}{130,30,30}
\definecolor{FdMGreen}{RGB}{25,100,60}

\setbeamercolor{structure}{fg=FdMBlue}
\setbeamercolor{title}{fg=white,bg=FdMBlue}
\setbeamercolor{frametitle}{fg=white,bg=FdMBlue}
\setbeamercolor{block title}{fg=white,bg=FdMBlue}
\setbeamercolor{block body}{bg=blue!3}
\setbeamercolor{alerted text}{fg=FdMRed}

\setmainfont{Latin Modern Roman}
\setsansfont{Latin Modern Sans}
\setmonofont{Latin Modern Mono}

\newcommand{\N}{\mathbb{N}}
\newcommand{\Z}{\mathbb{Z}}
\newcommand{\Q}{\mathbb{Q}}
\newcommand{\R}{\mathbb{R}}
\newcommand{\C}{\mathbb{C}}
\newcommand{\Pow}{\mathcal{P}}
\newcommand{\id}{\operatorname{id}}
\newcommand{\mcd}{\operatorname{mcd}}
\newcommand{\mcm}{\operatorname{mcm}}

\title[FdM I -- Clase 13]{Semigrupos, monoides y grupos}
\subtitle{Primeras estructuras algebraicas}
\author{Antonio Falcó Montesinos}
\institute{Universidad CEU Cardenal Herrera}
\date{Fundamentos de las Matemáticas I}

\begin{document}

\begin{frame}
  \titlepage
\end{frame}

\begin{frame}{Objetivos de la clase}
\begin{itemize}
  \item Definir semigrupos, monoides y grupos.
  \item Reconocer ejemplos básicos.
  \item Entender la noción de grupo abeliano.
\end{itemize}
\end{frame}

\begin{frame}{Mapa de la clase}
\tableofcontents
\end{frame}

\section{Estructuras con una operación}

\begin{frame}{Estructuras con una operación}
\begin{block}{Semigrupo}
\begin{itemize}
  \item Un semigrupo es un par \((A,*)\) con operación interna asociativa.
  \item Ejemplo: \((\mathbb{N},+)\).
  \item No se exige neutro ni inversos.
\end{itemize}
\end{block}
\end{frame}

\begin{frame}{Estructuras con una operación}
\begin{block}{Monoide}
\begin{itemize}
  \item Un monoide es un semigrupo con elemento neutro.
  \item Ejemplo: \((\mathbb{N},+,0)\).
  \item La notación debe indicar conjunto, operación y neutro.
\end{itemize}
\end{block}
\end{frame}

\begin{frame}{Estructuras con una operación}
\begin{block}{Grupo}
\begin{itemize}
  \item Un grupo tiene operación interna asociativa, neutro e inversos para todos sus elementos.
  \item Ejemplo: \((\mathbb{Z},+,0)\).
  \item No toda estructura con neutro es un grupo.
\end{itemize}
\end{block}
\end{frame}

\section{Lectura estructural}

\begin{frame}{Lectura estructural}
\begin{block}{Grupo abeliano}
\begin{itemize}
  \item Un grupo es abeliano si además la operación es conmutativa.
  \item \((\mathbb{Z},+,0)\) es abeliano.
  \item Las simetrías de una figura forman un grupo; no siempre es abeliano.
\end{itemize}
\end{block}
\end{frame}

\begin{frame}{Lectura estructural}
\begin{block}{Errores frecuentes}
\begin{itemize}
  \item Olvidar comprobar que la operación es interna.
  \item Confundir neutro con inverso.
  \item Suponer que toda operación conocida conserva sus propiedades en cualquier conjunto.
\end{itemize}
\end{block}
\end{frame}

\section{Profundización}

\begin{frame}{Profundización}
\begin{block}{Jerarquía}
\begin{itemize}
  \item Todo grupo es un monoide.
  \item Todo monoide es un semigrupo.
  \item Las implicaciones inversas no son ciertas en general.
\end{itemize}
\end{block}
\end{frame}

\begin{frame}{Profundización}
\begin{block}{Ejemplos comparados}
\begin{itemize}
  \item \((\mathbb{N},+)\) es monoide si \(0\in\mathbb{N}\).
  \item \((\mathbb{Z},+)\) es grupo abeliano.
  \item \((\mathbb{N},+)\) no es grupo porque \(1\) no tiene inverso aditivo en \(\mathbb{N}\).
\end{itemize}
\end{block}
\end{frame}

\begin{frame}{Profundización}
\begin{block}{Cómo probar que algo no es grupo}
\begin{itemize}
  \item Mostrar que la operación no es interna.
  \item Mostrar que no es asociativa.
  \item Mostrar que no existe neutro.
  \item Mostrar que algún elemento no tiene inverso.
\end{itemize}
\end{block}
\end{frame}

\begin{frame}{Cierre}
\begin{itemize}
  \item Semigrupo: asociatividad.
  \item Monoide: asociatividad y neutro.
  \item Grupo: asociatividad, neutro e inversos.
\end{itemize}
\end{frame}

\begin{frame}{Trabajo recomendado}
\begin{itemize}
  \item Releer las secciones correspondientes del manual.
  \item Identificar definiciones, símbolos nuevos y enunciados principales.
  \item Preparar dudas y ejemplos para el seminario asociado.
\end{itemize}
\end{frame}

\end{document}
